\begin{filecontents*}[overwrite]{verify_polynomial_inverse.py}
"""Exact checks accompanying the uniform polynomial-pullback proof."""
from fractions import Fraction
from pathlib import Path
import json


def add_term(poly, degree, coefficient):
    poly[degree] = poly.get(degree, Fraction(0)) + coefficient
    if not poly[degree]:
        del poly[degree]


def primitive(poly, degree):
    answer = {}
    for power, coefficient in poly.items():
        add_term(answer, power + degree,
                 coefficient * Fraction(degree, power + degree))
        add_term(answer, power + 1, coefficient / (power + 1))
    return answer


def differentiate(poly):
    return {power - 1: coefficient * power
            for power, coefficient in poly.items() if power}


def multiply_by_cover_derivative(poly, degree):
    answer = dict(poly)
    for power, coefficient in poly.items():
        add_term(answer, power + degree - 1, degree * coefficient)
    return answer


def serial(poly):
    return {str(power): str(poly[power]) for power in sorted(poly)}


checks = 0
families = []
cubic = []
for degree in range(3, 11):
    current = {1: Fraction(1)}
    for depth in range(1, 21):
        following = primitive(current, degree)
        assert differentiate(following) == multiply_by_cover_derivative(current, degree)
        assert max(following) == depth * degree + 1
        assert min(following) == depth + 1
        assert len(following) == depth + 1
        checks += 1
        if degree == 3 and depth <= 3:
            cubic.append(serial(following))
        current = following
    families.append({"degree": degree, "depth": 20,
                     "final_polynomial_degree": max(current),
                     "final_monomial_count": len(current)})

assert cubic == [
    {"2": "1/2", "4": "3/4"},
    {"3": "1/6", "5": "9/20", "7": "9/28"},
    {"4": "1/24", "6": "19/120", "8": "117/560", "10": "27/280"},
]

result = {"arithmetic": "fractions.Fraction (exact rational arithmetic)",
          "identities_checked": checks, "failures": 0,
          "families": families, "cubic_first_three_primitives": cubic,
          "scope": "Finite exact checks support, but do not replace, the all-degree/all-depth proof."}
target = Path(__file__).with_suffix(".json")
target.write_text(json.dumps(result, indent=2) + "\n")
print(json.dumps(result, indent=2))
\end{filecontents*}
\begin{filecontents*}[overwrite]{verify_cover_classifier.py}
"""Exact finite classification checks on polynomial covers (SymPy).

Acceptance returns a polynomial q with D^N(q)=f. Rejection returns an
exact linear functional separating f from the bounded polynomial image.
The infinite-order theorem and the degree bound are proved in the paper.
"""
from pathlib import Path
import json
import sympy as S

t=S.symbols('t')

def delta(f,p):
    return S.cancel(S.diff(f,t)/S.diff(p,t))

def iterate(f,p,n):
    for _ in range(n):
        f=delta(f,p)
    return f

def classify(name,p,f,expected):
    p,f=S.expand(p),S.cancel(f)
    num,den=S.fraction(f)
    if f == 0:
        return dict(name=name,stable=True,N=0,degree_bound=0,witness='0')
    dp=S.Poly(S.diff(p,t),t,domain=S.QQ)
    N=0
    pole_data=[]
    for fac,m in S.factor_list(den,t)[1]:
        factor=S.Poly(fac,t,domain=S.QQ)
        val=0
        z=dp
        while z.rem(factor).is_zero:
            val+=1
            z=z.exquo(factor)
        e=val+1
        N=max(N,(m+e-1)//e)
        pole_data.append(dict(factor=str(fac),pole_order=int(m),ramification_index=int(e)))
    d=S.degree(p,t)
    beta=S.degree(num,t)-S.degree(den,t)
    B=int(max(d*N+beta,d*(N-1)))
    assert B>=0
    images=[iterate(t**j,p,N) for j in range(B+1)]
    common=S.Poly(1,t,domain=S.QQ)
    for g in images+[f]:
        common=S.lcm(common,S.Poly(S.fraction(g)[1],t,domain=S.QQ))
    cols=[S.Poly(S.cancel(g*common.as_expr()),t,domain=S.QQ) for g in images]
    target=S.Poly(S.cancel(f*common.as_expr()),t,domain=S.QQ)
    max_degree=max([0,int(target.degree())]+[int(q.degree()) for q in cols if not q.is_zero])
    A=S.Matrix([[q.nth(i) for q in cols] for i in range(max_degree+1)])
    b=S.Matrix([target.nth(i) for i in range(max_degree+1)])
    rank=A.rank()
    consistent=rank==A.row_join(b).rank()
    assert consistent==expected,name
    data=dict(name=name,cover=str(p),f=str(f),N=int(N),degree_bound=B,
              poles=pole_data,matrix_shape=list(A.shape),stable=consistent)
    if consistent:
        sol,params=A.gauss_jordan_solve(b)
        sol=sol.subs({z:0 for z in params})
        q=S.expand(sum(sol[j]*t**j for j in range(B+1)))
        assert S.cancel(iterate(q,p,N)-f)==0
        data['polynomial_witness']=str(q)
    else:
        w=next(v for v in A.T.nullspace() if (v.T*b)[0]!=0)
        assert w.T*A==S.zeros(1,A.cols)
        assert (w.T*b)[0]!=0
        data['separating_functional']=[str(v) for v in w]
        data['nonzero_pairing']=str((w.T*b)[0])
    data['verified']=True
    print(json.dumps(data),flush=True)
    return data

if __name__=='__main__':
    p=t**3+t
    cases=[
        ('cubic inverse',p,t,True),
        ('first derivative',p,1/(3*t*t+1),True),
        ('weighted first derivative',p,t/(3*t*t+1),True),
        ('third derivative of a polynomial',p,iterate(t*t,p,3),True),
        ('simple ordinary pole',p,1/t,False),
        ('double critical denominator',p,1/(3*t*t+1)**2,False),
        ('quartic inverse second derivative',t**4+t,iterate(t**3,t**4+t,2),True),
        ('quintic inverse',t**5+t,t,True),
        ('quintic ordinary pole',t**5+t,1/t,False),
        ('cyclic fractional power',t**4,t**-1,True),
        ('cyclic logarithmic obstruction',t**4,t**-4,False),
        ('cyclic delayed logarithmic obstruction',t**4,t**-8,False),
        ('cyclic negative fractional power',t**4,t**-7,True),
        ('multiple ramification', (t*t+1)**3-1,
         iterate(t,(t*t+1)**3-1,3),True),
    ]
    results=[classify(*case) for case in cases]
    Path(__file__).with_suffix('.json').write_text(json.dumps(results,indent=2)+'\n')
\end{filecontents*}
\documentclass[11pt,a4paper]{article}
\usepackage[margin=25mm]{geometry}
\usepackage[T1]{fontenc}
\usepackage{lmodern,amsmath,amssymb,amsthm,mathtools,booktabs,array}
\usepackage{xcolor,microtype,fancyhdr}
\usepackage[colorlinks=true,linkcolor=blue!45!black,urlcolor=blue!45!black,citecolor=blue!45!black]{hyperref}
\hypersetup{pdftitle={Polynomial Inverses and Algebraic Integration: Counterexamples to Chen's Stability Conjecture},pdfauthor={Keio}}
\newtheorem{theorem}{Theorem}[section]
\newtheorem{lemma}[theorem]{Lemma}
\newtheorem{proposition}[theorem]{Proposition}
\newtheorem{corollary}[theorem]{Corollary}
\theoremstyle{remark}\newtheorem{remark}[theorem]{Remark}
\setlength{\emergencystretch}{3em}
\setlength{\headheight}{14pt}
\pagestyle{fancy}\fancyhf{}
\fancyhead[L]{\small Keio}
\fancyhead[R]{\small Polynomial inverses and algebraic integration}
\fancyfoot[C]{\thepage}
\title{Polynomial Inverses and Algebraic Integration\\[5pt]
\large Counterexamples to Chen's Algebraic-Stability Conjecture\\
and a Complete Criterion on Fixed Polynomial Covers}
\author{Keio\thanks{Research development, independent internal checks, and manuscript preparation with ChatGPT.}}
\date{2 November 2026}
\begin{document}
\maketitle
\begin{abstract}
Chen's 2022 algebraic-stability conjecture proposes that algebraic functions
admitting arbitrarily many algebraic primitives are precisely finite
polynomial linear combinations of rational powers of shifted variables.
We give counterexamples in every degree $d\ge3$: an inverse branch of
$t^d+t=x$ has an infinite primitive tower in $\mathbb C[t]$, yet generates
a non-Galois extension of $\mathbb C(x)$ and therefore cannot belong to
the proposed class. More generally, every polynomial inverse is
algebraically stable, and it lies in that class exactly when its defining
polynomial has the form $a(t-b)^d+c$. For a fixed polynomial cover
$x=P(t)$, we classify all algebraically stable elements of $\mathbb C(t)$
as $\bigcup_{n\ge0}\delta^n\mathbb C[t]$, where
$\delta=P'(t)^{-1}d/dt$. Finite pole orders give an explicit bound reducing
membership to one polynomial linear system. The proofs use classical
substitution, field trace, valuations, and Galois theory. They refute
the stated 2022 conjecture, but do not solve the full elementary-function
classification in Chen--Kauers Problem~8. Historical priority is not
certified. Two exact verification programs are included in the single
TeX source.
\end{abstract}

\section{The historical target and the scope of the result}
Original Problem~8 of Chen--Kauers~\cite{CK2016} asks which elementary
functions remain elementary after arbitrarily many integrations.
This is a broad classification problem. In his later study,
Chen~\cite[Section~5]{Chen2022} proposed a more specific description
for functions whose repeated primitives can all be chosen algebraic.
The counterexamples below concern that exact description. They are
not a redefinition of the original elementary-function problem.

We distinguish two properties. A function is \emph{algebraically stable}
if it admits a coherent infinite sequence of primitives in
$\overline{\mathbb C(x)}$. It is \emph{elementary-stable} if those primitives
may use elementary extensions, including logarithms. Algebraic stability
implies elementary stability; the converse fails even on rationally
parametrized algebraic curves. The distinction also separates both
notions from stability of the minimum order of a differential equation,
studied in~\cite{CFGL2023}.

\paragraph{Source verification and priority.}
The conjecture used here occurs in the conclusion of arXiv:2202.06305v1,
printed page~14, and is retained in the author's eight-page
conference-format manuscript, Section~5, PDF page~7. Its coefficient
and exponent restrictions are reproduced mathematically below.
A 2025 paper continues to study the related algebraic-stability problem
through special fields~\cite[Proposition~5.3]{SSS2025}. A search during
preparation did not identify an earlier identical disproof. This limited
search does not establish firstness, and the methods themselves are
classical. The claims are the displayed theorems and the resulting
contradiction to the stated conjecture, not a new foundational principle
or an independently certified priority claim.
% Manuscript fragment: requires amsmath, amssymb, amsthm and theorem/proposition
% environments. Bibliography key Chen2022 should refer to Stability Problems
% in Symbolic Integration, ISSAC 2022, DOI 10.1145/3476446.3535502.
\section{Polynomial inverses and Chen's algebraic-stability conjecture}
\label{sec:polynomial-inverses}

We address the specific algebraic-stability conjecture proposed in Section~5
of Chen's \emph{Stability Problems in Symbolic Integration}~\cite{Chen2022}. Its asserted
description of the elements stable under $d/dx$ in
$\overline{\mathbb C(x)}$ is the class of finite sums
\begin{equation}\label{eq:chen-proposed-algebraic-class}
 \sum_{i=1}^{r}p_i(x)(x-c_i)^{\alpha_i},\qquad
 p_i\in\mathbb C[x],\quad c_i\in\mathbb C,\quad
 \alpha_i\in\mathbb Q\setminus\{-1,-2,\ldots\}.
\end{equation}
Here stability means the existence of a sequence $f_0,f_1,\ldots$ in the
algebraic closure with $f_0=f$ and $Df_{n+1}=f_n$ for every $n\ge0$.
We give counterexamples to the proposed necessity in
\eqref{eq:chen-proposed-algebraic-class}. This is a statement about that
specific conjecture; it is not a classification of all elementary-stable
functions and does not close Chen--Kauers Problem~8. No assertion of historical
priority is needed for the mathematical arguments below.

\begin{proposition}[Polynomial pullback primitives]
\label{prop:polynomial-pullback}
Let $P\in\mathbb C[t]$ be nonconstant and identify $x=P(t)$. Every
$Q(t)\in\mathbb C[t]$ is stable under $D=d/dx$ in the finite algebraic
extension $\mathbb C(t)/\mathbb C(x)$. In particular, every local inverse
branch $t=t(x)$ of a polynomial admits an infinite tower of algebraic
primitives in the same function field.
\end{proposition}

\begin{proof}
The derivation extends uniquely to the function field, where
\[
 D=\frac1{P'(t)}\frac d{dt}.
\]
Define a linear map on polynomials by
\begin{equation}\label{eq:pullback-primitive-map}
 (J_PQ)(t)=\int_0^t P'(u)Q(u)\,du.
\end{equation}
The integral in \eqref{eq:pullback-primitive-map} is ordinary polynomial
antidifferentiation, so $J_PQ\in\mathbb C[t]$, and the chain rule gives
$D(J_PQ)=Q$. Consequently $f_n=J_P^nQ$ is the required infinite tower.
Although $D$ need not preserve the polynomial subring, the right-inverse
construction does; no assertion that $\mathbb C[t]$ is a differential subring
is being made.
\end{proof}

\begin{proposition}[An abelian-extension obstruction]
\label{prop:abelian-obstruction}
Every function of the form \eqref{eq:chen-proposed-algebraic-class} generates
a Galois extension of $\mathbb C(x)$ whose Galois group is abelian.
\end{proposition}

\begin{proof}
Write $\alpha_i=a_i/m_i$ with $a_i\in\mathbb Z$ and $m_i\ge1$, and choose
$u_i$ with $u_i^{m_i}=x-c_i$. Because $\mathbb C$ contains all roots of unity,
$\mathbb C(x,u_i)/\mathbb C(x)$ is a cyclic Galois extension, allowing the
trivial case. Its splitting polynomial is $T^{m_i}-(x-c_i)$. The compositum
\[
 E=\mathbb C(x,u_1,\ldots,u_r)
\]
is finite Galois, and its Galois group embeds into a product of cyclic groups,
so is abelian. The displayed sum belongs to $E$. Every subgroup of an abelian
group is normal; therefore every intermediate field, including the field
generated by the sum, is Galois over $\mathbb C(x)$, with abelian group.
\end{proof}

\begin{theorem}[Which polynomial inverses have the proposed form]
\label{thm:polynomial-inverse-classification}
Let $P\in\mathbb C[t]$ have degree $d\ge1$, and identify $x=P(t)$. The
function $t$ has a representation of the form
\eqref{eq:chen-proposed-algebraic-class} if and only if
\begin{equation}\label{eq:affine-power-polynomial}
 P(t)=a(t-b)^d+c
 \qquad(a\in\mathbb C^*,\ b,c\in\mathbb C).
\end{equation}
\end{theorem}

\begin{proof}
The sufficiency follows from
$t=b+a^{-1/d}(x-c)^{1/d}$. Conversely, a representation of the proposed form
implies by Proposition~\ref{prop:abelian-obstruction} that
$\mathbb C(t)/\mathbb C(P(t))$ is Galois. It has degree $d$, so its
automorphism group $G$ has order $d$.

For $\sigma\in G$, write $\sigma(t)=r_\sigma(t)\in\mathbb C(t)$. The identity
$P(r_\sigma(t))=P(t)$ shows that $r_\sigma$ has no finite pole and that
$\deg r_\sigma=1$. Hence the maps $r_\sigma$ form a finite group of affine
transformations. The point
\[
 b=\frac1{|G|}\sum_{\sigma\in G}r_\sigma(0)
\]
is fixed by every map $r_\sigma$: applying any affine group element commutes
with this average and permutes the terms. In the coordinate $u=t-b$, the
group therefore consists of multiplications $u\mapsto\lambda u$. It is a
subgroup of $\mathbb C^*$ of order $d$, hence the group of $d$th roots of
unity. Invariance of the degree-$d$ polynomial $P(b+u)$ then forces all its
nonconstant monomials to have degrees divisible by $d$. Consequently
$P(b+u)=au^d+c$ with $a\ne0$, proving \eqref{eq:affine-power-polynomial}.
\end{proof}

\begin{theorem}[Counterexamples in every degree at least three]
\label{thm:polynomial-inverse-counterexamples}
For any integer $d\ge3$, let $t_d$ satisfy
\begin{equation}\label{eq:inverse-family}
 t_d^d+t_d=x.
\end{equation}
Then $t_d$ is stable under $d/dx$ in $\overline{\mathbb C(x)}$, but it has
no representation of the form \eqref{eq:chen-proposed-algebraic-class}.
Hence the proposed necessity in Chen's algebraic-stability conjecture is false.
\end{theorem}

\begin{proof}
Stability follows from Proposition~\ref{prop:polynomial-pullback} with
$P(t)=t^d+t$ and $Q(t)=t$. It remains to prove that
$\mathbb C(t)/\mathbb C(t^d+t)$ is not Galois.

The polynomial $T^d+T-x$ is irreducible over $\mathbb C(x)$. Indeed
$x-T^d-T$ is irreducible in $\mathbb C[x,T]$: viewed as a polynomial in $x$
over $\mathbb C[T]$, its leading coefficient is a unit and its degree is one.
Gauss's lemma then proves the asserted irreducibility over $\mathbb C(x)$.
Thus the extension has degree $d$ and is separable.

Suppose an automorphism over $\mathbb C(t^d+t)$ sends $t$ to $r(t)\in
\mathbb C(t)$. Then
\begin{equation}\label{eq:deck-equation}
 r(t)^d+r(t)=t^d+t.
\end{equation}
A finite pole of $r$ would give a finite pole on the left, with the highest
pole order coming from $r^d$ and therefore unable to cancel. The right side
has no finite poles, so $r$ is a polynomial. Comparing degrees in
\eqref{eq:deck-equation} gives $\deg r=1$; write $r(t)=at+b$ with $a\ne0$.
Since $d\ge3$, comparison of the coefficient of $t^{d-1}$ gives
$da^{d-1}b=0$, hence $b=0$. Leading coefficients give $a^d=1$, while the
coefficient of $t$ gives $a=1$. Thus the identity is the only automorphism.
An extension of degree $d>1$ with only one automorphism is not Galois.

If $t_d$ had the proposed representation, Proposition~\ref{prop:abelian-obstruction}
would make its generated extension Galois, a contradiction.
\end{proof}

The same obstruction excludes even finite sums of rational multiples of
rational-function radicals whose compositum is abelian. The proof uses
polynomial substitution and elementary Galois theory; it does not introduce
a new foundational principle or rely on a conjectural classification.

\paragraph{The cubic example and its first three primitives.}
For $x=t^3+t$, choose all integration constants to make the polynomials vanish
at $t=0$. The first three primitives of $t$ are
\begin{align*}
 F_1(t)&=\frac34t^4+\frac12t^2,\\
 F_2(t)&=\frac9{28}t^7+\frac9{20}t^5+\frac16t^3,\\
 F_3(t)&=\frac{27}{280}t^{10}+\frac{117}{560}t^8
          +\frac{19}{120}t^6+\frac1{24}t^4.
\end{align*}
They satisfy $DF_1=t$, $DF_2=F_1$, and $DF_3=F_2$, with
$D=(3t^2+1)^{-1}d/dt$. An independent check of non-Galoisness in degree three
is the discriminant
\[
 \operatorname{disc}_T(T^3+T-x)=-4-27x^2,
\]
which has two simple zeros and is not a square in $\mathbb C(x)$.

More generally the tower for \eqref{eq:inverse-family} can be computed using
the monomial rule
\[
 J_P(t^m)=\frac{d}{m+d}t^{m+d}+\frac1{m+1}t^{m+1}.
\]
Starting from $t$, its $n$th member has degree $nd+1$. This exact recurrence
also gives reproducible symbolic verification of as many initial levels as
desired; such checks illustrate the uniform proof rather than replace it.
\section{Complete algebraic-stability classification on a polynomial cover}
\label{sec:classification}
Fix a nonconstant $P\in\mathbb C[t]$, put $x=P(t)$, and write
$K=\mathbb C(t)$ and $\delta=P'(t)^{-1}d/dt$.
The next result classifies all rational functions in this presentation,
with algebraic primitives allowed in the entire algebraic closure.
It is more restrictive than elementary integration, where logarithms
are permitted, and more general than considering only the inverse
function $t$ itself.

\begin{lemma}[Descent of algebraic primitives]\label{lem:trace}
Let $K$ be a characteristic-zero differential field. If $f\in K$ and
$\delta^n g=f$ for an element $g$ algebraic over $K$, then there is
$h\in K$ with $\delta^n h=f$.
\end{lemma}
\begin{proof}
The finite extension $L=K(g)$ is stable under the uniquely extended
derivation: differentiating the minimal polynomial expresses $\delta g$
in $L$, since its derivative with respect to its indeterminate is nonzero
at $g$. Every $K$-embedding of $L$ into an algebraic closure commutes
with the extended derivation, by uniqueness. Thus field trace commutes
with $\delta$ and its powers. The normalized trace
$h=[L:K]^{-1}\operatorname{Tr}_{L/K}(g)$ satisfies $\delta^n h=f$.
\end{proof}

\begin{lemma}[Finite pole growth]\label{lem:pole-growth}
At a finite point $a$, put
$e_a=\operatorname{ord}_{t=a}(P(t)-P(a))\ge1$.
If $g\in K$ has a pole of order $k>0$ at $a$, then $\delta^n g$
has a pole of exactly $k+ne_a$ at $a$, for every $n\ge0$.
\end{lemma}
\begin{proof}
With $u=t-a$, write $P(t)-P(a)=c u^{e_a}+O(u^{e_a+1})$ and
$g=b u^{-k}+O(u^{1-k})$, with $bc\ne0$. Then
\[
 \delta g=-\frac{kb}{e_a c}u^{-k-e_a}
                 +O(u^{1-k-e_a}).
\]
The leading coefficient is nonzero in characteristic zero. Iterating
the same calculation proves the assertion.
\end{proof}

\begin{theorem}[Classification and a finite decision bound]
\label{thm:cover-classification}
For nonzero $f\in\mathbb C(t)$, let $m_a$ be its pole order at each
finite pole $a$. Define
\begin{equation}\label{eq:N-bound}
 N(f)=\max_{a\text{ a finite pole of }f}
              \left\lceil\frac{m_a}{e_a}\right\rceil,
 \qquad \max\varnothing=0.
\end{equation}
The following are equivalent:
\begin{enumerate}
\item $f$ is stable under $d/dx$ in $\overline{\mathbb C(x)}$;
\item $f=\delta^Nq$ for some $q\in\mathbb C[t]$, with $N=N(f)$;
\item $f$ has an algebraic $N(f)$-fold primitive.
\end{enumerate}
In particular,
\begin{equation}\label{eq:stable-union}
 \operatorname{Stab}(\delta,\overline{\mathbb C(x)})\cap\mathbb C(t)
   =\bigcup_{n\ge0}\delta^n\mathbb C[t].
\end{equation}
For $f=0$ all assertions hold trivially.
\end{theorem}
\begin{proof}
The implication (1)$\Rightarrow$(3) is immediate. For (3)$\Rightarrow$(2),
Lemma~\ref{lem:trace} descends an algebraic $N$-fold primitive to some
$q\in K$. Suppose $q$ has a finite pole at $a$, of order $k>0$.
Lemma~\ref{lem:pole-growth} implies that $f$ has a pole there of order
$m_a=k+Ne_a>Ne_a$. But the choice of $N$ gives $Ne_a\ge m_a$, a
contradiction. A pole of $q$ outside the pole set of $f$ is excluded by
the same lemma. Hence $q$ has no finite poles and is a polynomial.
If $N=0$, the hypothesis says simply that $f$ has no finite poles, with
the same conclusion.

For (2)$\Rightarrow$(1), polynomial primitives are supplied by $J_P$
from~\eqref{eq:pullback-primitive-map}. A coherent primitive sequence
for $f=\delta^Nq$ begins with
\[
 \delta^Nq,\ \delta^{N-1}q,\ldots,q,\ J_Pq,\ J_P^2q,\ldots.
\]
Every element lies in $K$. This proves the equivalence and the union
description. The union is increasing: $\mathbb C[t]\subseteq
\delta\mathbb C[t]$ because $\delta J_P$ is the identity on polynomials.
\end{proof}

\begin{proposition}[A polynomial degree bound]\label{prop:q-degree}
Let $d=\deg P$ and, for a reduced fraction $f=A/B\ne0$, put
$\beta=\deg A-\deg B$. If $f=\delta^Nq$ for a polynomial $q$, then
\begin{equation}\label{eq:q-degree}
 \deg q\le B_f:=\max\{Nd+\beta,\ d(N-1)\}.
\end{equation}
Here $\beta$ is the growth degree at infinity, the negative of the
valuation at infinity. For $N=0$, the bound is $\deg q=\beta$.
\end{proposition}
\begin{proof}
Write $P(t)=c t^d+O(t^{d-1})$. If $M=\deg q>d(N-1)$, then successive
leading terms under $\delta$ have degrees $M-d,M-2d,\ldots,M-Nd$,
and their combined leading coefficient contains the nonzero factor
\[
 (dc)^{-N}\prod_{j=0}^{N-1}(M-jd).
\]
Consequently $\beta=M-Nd$ and $M=Nd+\beta$. If
$M\le d(N-1)$ there is nothing to prove. The case $N=0$ is immediate.
\end{proof}

Theorems~\ref{thm:cover-classification} and~\ref{prop:q-degree} give a
terminating decision procedure over an effective coefficient field.
Factor the reduced denominator of $f$, calculate its pole orders and
the corresponding multiplicities in $P'$, and form $N(f)$.
Then substitute
\[
 q(t)=\sum_{j=0}^{B_f}c_jt^j
\]
into $\delta^Nq=f$, clear denominators, and solve the finite linear
system. If $B_f<0$, no nonzero polynomial candidate exists, so the
answer is negative. A solution provides the whole primitive tower;
inconsistency proves that no algebraic $N$-fold primitive, and hence
no algebraic stability, is possible. Over $\mathbb Q$, factorization
over $\mathbb Q$ suffices for computing $N$: all roots of a given
irreducible factor have the same denominator and $P'$ multiplicities.

This is a complete classification in the stated fixed polynomial
cover. It does not assert that every algebraic function belongs to
such a presentation.

\section{Elementary stability and the scope of Problem 8}
\label{sec:elementary}
The distinction between algebraic and elementary primitives is essential.
An algebraic function may fail the classification above and nevertheless
admit arbitrarily many elementary primitives.

\begin{lemma}[Polynomial moments and repeated integration]
\label{lem:moments}
In a differential field with $Dx=1$, if $H_j'=x^jf$ for all $j\ge0$,
then a coherent primitive sequence is given by $G_0=f$ and
\begin{equation}\label{eq:moment-primitive}
 G_n=\frac1{(n-1)!}\sum_{j=0}^{n-1}
       (-1)^j\binom{n-1}{j}x^{n-1-j}H_j,\qquad n\ge1.
\end{equation}
Conversely, the existence of all repeated primitives implies the
integrability of all $x^jf$ in the same field.
\end{lemma}
\begin{proof}
Differentiating~\eqref{eq:moment-primitive}, the terms containing $f$
cancel by the binomial theorem when $n>1$, and the remaining terms give
$G_{n-1}$. For $n=1$, the derivative is $f$. Conversely, if
$D^{j+1}g=f$, repeated integration by parts expresses a primitive of
$x^jf$ as a finite sum of polynomial multiples of derivatives of $g$.
\end{proof}
This equivalence is already used in Chen--Kauers~\cite{CK2016} and
Chen~\cite[Lemma~2.10]{Chen2022}; it is recalled to make the scope clear.

\begin{proposition}[Rationally parametrized elementary stability]
\label{prop:rational-cover}
Let $x=R(t)$ for any nonconstant rational function
$R\in\mathbb C(t)$. Every $f\in\mathbb C(t)$ is elementary-stable
over $\mathbb C(x)$.
\end{proposition}
\begin{proof}
For every $j\ge0$, substitution gives
\[
 x^jf\,dx=R(t)^j f(t)R'(t)\,dt,
\]
a rational differential in $t$. Its primitive is a rational function
plus a finite constant linear combination of logarithms of rational
functions in $t$, by partial fractions. Since $t$ is algebraic over
$\mathbb C(x)$, these primitives are elementary over that field.
Lemma~\ref{lem:moments} constructs all repeated elementary primitives.
\end{proof}

For example, on $x=t^3+t$, the function $f=1/t$ has
$N(f)=1$ but no algebraic primitive:
\[
 \int f\,dx=\int\frac{3t^2+1}{t}\,dt
          =\frac32t^2+\log t.
\]
Its nonzero simple residue prevents a rational primitive in $\mathbb C(t)$,
and Lemma~\ref{lem:trace} prevents an algebraic primitive in any finite
extension. Nevertheless it is elementary-stable by
Proposition~\ref{prop:rational-cover}.

Chen--Kauers Problem~8 asks for the class of \emph{all} elementary-stable
functions. The results here settle neither higher-genus algebraic
presentations in general nor arbitrary nested exponential and logarithmic
extensions. They refute a particular proposed algebraic classification
and give exact positive and negative results on a substantial family.
The full original problem remains unresolved by this manuscript.

\section{Concrete examples and exact verification}
\subsection{An elementary real branch}
The cubic counterexample has a unique real branch, because
$t\mapsto t^3+t$ is strictly increasing. It can be written explicitly as
\begin{equation}\label{eq:real-cubic}
 t(x)=\frac2{\sqrt3}\sinh\!\left(
       \frac13\operatorname{arsinh}\!\left(\frac{3\sqrt3}{2}x\right)
       \right).
\end{equation}
The identity $\sinh(3u)=4\sinh^3u+3\sinh u$ verifies $t^3+t=x$.
Thus the example can also be presented using familiar elementary
functions. Its exclusion from Chen's proposed class does not mean it
is inexpressible using nested radicals or elementary compositions;
it excludes the particular finite sum in
\eqref{eq:chen-proposed-algebraic-class}.

Its first primitive is already enough to display the mechanism:
\[
 \frac{d}{dx}\left(\frac34t(x)^4+\frac12t(x)^2\right)
 =\frac{3t^3+t}{3t^2+1}=t.
\]
The map $J_P$ supplies every later primitive without introducing a
larger field or an unproved limiting step.

\subsection{Positive and negative classification examples}
For $x=t^3+t$, set $q(t)=3t^2+1$. Representative outcomes of the finite
algorithm are shown below. In this table ``stable'' means
\emph{algebraically} stable. All listed functions are elementary-stable
by Proposition~\ref{prop:rational-cover}.
\begin{center}
\begin{tabular}{@{}lrrl@{}}
\toprule
Input & $N(f)$ & $B_f$ & Algebraic outcome\\
\midrule
$t$ & $0$ & $1$ & Stable: polynomial tower\\
$1/q$ & $1$ & $1$ & Stable: $\delta(t)=1/q$\\
$t/q$ & $1$ & $2$ & Stable: $\delta(t^2/2)=t/q$\\
$\delta^3(t^2)$ & $3$ & $6$ & Stable: witness $t^2$\\
$1/t$ & $1$ & $2$ & No algebraic first primitive\\
$1/q^2$ & $1$ & $0$ & No algebraic first primitive\\
\bottomrule
\end{tabular}
\end{center}
The second negative case follows because a candidate polynomial
primitive would have degree at most zero and therefore derivative zero.
This conclusion concerns primitives in the entire algebraic closure,
by trace descent, not just a failed choice of rational ansatz.

On the cyclic cover $x=t^4$, the same algorithm accepts $t^{-1}$ and
$t^{-7}$ but rejects $t^{-4}$ and $t^{-8}$. For example,
\[
 \delta^2\left(-\frac{16}{3}t\right)=t^{-7},\qquad
 t^{-8}=x^{-2}.
\]
The latter has a first algebraic primitive $-x^{-1}$, followed by a
logarithmic obstruction. This tests the distinction between one-time
algebraic integrability and stability.

\subsection{Verification record}
The first embedded program uses only Python's exact rational arithmetic.
It verifies the primitive recurrence for each $d=3,\ldots,10$ through
depth $20$: all $160$ identities pass. It also checks the first three
cubic primitives printed above and the degree of every tested tower
element. These finite checks supplement the uniform proof, which
applies to every $d\ge3$ and every depth.

The second program implements the classification in
Theorem~\ref{thm:cover-classification} over $\mathbb Q$. It checks $14$
examples, including ordinary poles, simple and multiple ramification,
and delayed logarithmic obstructions. Every positive case returns a
polynomial $q$ and verifies $\delta^Nq=f$ exactly. Every negative case
returns a row vector $w$ for the cleared linear system $Ac=b$ such that
\[
 wA=0,\qquad wb\ne0.
\]
This is an exact inconsistency certificate. All $14$ outcomes pass their
independent identity or separating-functional checks. The proof of the
degree bound is essential: an inconsistent arbitrary truncated ansatz
would not certify nonexistence.

\paragraph{Reproduction from a single source.}
The source includes the bibliography and both Python programs in
\texttt{filecontents*} blocks. Normal LaTeX compilation extracts the
programs without executing them. Run
\begin{verbatim}
python3 verify_polynomial_inverse.py
python3 verify_cover_classifier.py
\end{verbatim}
The first program uses the standard Python library; the second requires
SymPy. Both produce JSON records. No shell escape, external bibliography,
or separate figure file is required to compile the paper.

\section{Consequences and remaining questions}
The conjectured algebraic classification fails by explicit counterexamples
of every degree at least three. The obstruction is structural: the
proposed expression class lies in abelian extensions, whereas polynomial
inverse branches need not generate normal extensions at all. Repeated
algebraic integration is therefore compatible with this nonabelian
field-theoretic behavior.

For polynomial covers, the corrected description
\eqref{eq:stable-union} is complete and effective. It explains both the
counterexamples and their stable rational derivatives, and gives finite
negative certificates for functions that acquire logarithms after a
bounded number of integrations. It does not propose a universal
replacement classification for all algebraic functions.

The full Chen--Kauers Problem~8 still requires arbitrary elementary
towers to be treated. This manuscript does not prove that every
elementary-stable function is D-finite, does not classify positive-genus
algebraic function fields in general, and does not certify the absence
of earlier equivalent results. The counterexample proof is unconditional
and self-contained; external review and a broader priority assessment
remain separate tasks.

\begin{thebibliography}{9}
\bibitem{CK2016}
S. Chen and M. Kauers,
\emph{Some open problems related to creative telescoping},
arXiv:1609.03768, 2016, original Problem~8.
\url{https://arxiv.org/abs/1609.03768}.

\bibitem{Chen2022}
S. Chen, \emph{Stability problems in symbolic integration},
Proceedings of ISSAC 2022, pp.~517--524,
\url{https://doi.org/10.1145/3476446.3535502}.
Preprint: \url{https://arxiv.org/abs/2202.06305}.
The conjecture is in Section~5, printed p.~14 of the preprint.
Author's conference-format manuscript:
\url{http://www.mmrc.iss.ac.cn/~schen/Chen2022stable.pdf}.

\bibitem{CFGL2023}
S. Chen, R. Feng, Z. Guo, and W. Lu,
\emph{Stability problems on D-finite functions},
Proceedings of ISSAC 2023, pp.~164--172,
\url{https://doi.org/10.1145/3597066.3597085}.
\url{https://arxiv.org/abs/2311.05897}.

\bibitem{SSS2025}
C. Sahu, M. Seiss, and V. R. Srinivasan,
\emph{Iterated and generalized iterated integrals},
arXiv:2504.04846, 2025;
Journal of Algebra and Its Applications,
\url{https://doi.org/10.1142/S0219498826502865}.
\url{https://arxiv.org/abs/2504.04846}.
\end{thebibliography}
\end{document}
